\section{Sums, number systems, inequalities, and square roots}\label{ch01:sec:notation}
To name the numbers in a list, we often use one letter with a small label below and to the right, called a \emph{subscript}: $a_1,a_2,a_3,a_4$ are the first four numbers of a list. A subscript is a label, not a multiplication: $a_2$ is the second number, whereas $2a$ is twice $a$. The notation
\[
\sum_{j=1}^{4}a_j=a_1+a_2+a_3+a_4
\]
means: add $a_1$, $a_2$, $a_3$, and $a_4$. The symbol $\sum$ is the Greek capital letter sigma, an S for ``sum.'' The letter $j$ is an \emph{index}, a counter that runs from the value below the sigma to the value above it.

If $a_1=1$, $a_2=3$, $a_3=5$, and $a_4=7$, the sum is $1+3+5+7=16$. Renaming the index does not change the sum, because the index is only a counter: $\sum_{r=1}^{4}a_r$ is also $16$. Changing the value above the sigma from $4$ to $3$ does: $\sum_{j=1}^{3}a_j=1+3+5=9$.

These totals are \emph{perfect squares}, squares of integers: $1+3+5=3^2$ and $1+3+5+7=4^2$, and likewise $1=1^2$ and $1+3=2^2$. Is the sum of the first $n$ odd positive integers always $n^2$? As in Section~\ref{ch01:sec:pattern}, examples suggest a pattern, but only a proof can settle it; Theorem~\ref{thm:odds} in Chapter~\ref{ch:induction} supplies one.

The quantity added can be any expression involving the index; substitute each value of the index in turn and add. With the same numbers,
\[
\sum_{j=1}^{3}(a_j+5)=(a_1+5)+(a_2+5)+(a_3+5)=6+8+10=24,
\]
and the index itself may appear in the expression:
\[
\sum_{j=1}^{3}j^2=1^2+2^2+3^2=1+4+9=14.
\]

A \emph{set} is a collection of objects, called its \emph{members}; curly braces list the members, and Chapter~\ref{ch:maps} studies sets carefully. Several sets of numbers have standard names.
\begin{itemize}
\item $\NN=\{0,1,2,3,\ldots\}$ is the set of \emph{natural numbers}. In this book zero counts as a natural number; when we want only the positive integers $1,2,3,\ldots$, we say so.
\item $\ZZ=\{\ldots,-2,-1,0,1,2,\ldots\}$ is the set of integers.
\item $\QQ$ is the set of \emph{rational numbers}: numbers that can be written as a fraction $a/b$ in which $a$ and $b$ are integers and $b\ne0$, such as $3/4$, $-5/2$, and $7=7/1$.
\item $\RR$ is the set of \emph{real numbers}: all the numbers on the number line, including some that are not fractions of integers. By Pythagoras' theorem, the diagonal of a square with sides of length $1$ has a length whose square is $1^2+1^2=2$; that length is $\sqrt2$, and Theorem~\ref{thm:sqrt2} in Chapter~\ref{ch:proof} proves that $\sqrt2$ is not such a fraction.
\end{itemize}
Each of these sets contains the one before it.

Four signs compare real numbers: $x<y$ means that $x$ is strictly smaller than $y$, and $x\le y$ means that $x$ is smaller than or equal to $y$, that is, at most $y$; likewise $x>y$ means strictly larger, and $x\ge y$ means at least $y$. So $3\le3$ is true, but $3<3$ is false. The signs $<$ and $>$ give \emph{strict} inequalities, and $\le$ and $\ge$ \emph{non-strict} ones. The two kinds disagree only when the two sides are equal, and many mistakes in proofs happen at that boundary case. A chain such as $0\le x\le1$ states two conditions at once: $0\le x$ and $x\le1$.

A number $x$ is \emph{positive} if $x>0$, \emph{negative} if $x<0$, and \emph{nonnegative} if $x\ge0$. So zero is nonnegative, but neither positive nor negative.

A number whose square is $t$ is called a \emph{square root} of $t$. We accept without proof that every nonnegative real number $t$ has exactly one nonnegative square root, written $\sqrt t$; this is the first of the starting facts listed in \emph{To the reader}. For example, $\sqrt9=3$: both $3$ and $-3$ have square $9$, and the square-root sign picks the nonnegative one. A negative number has no real square root, because the square of a real number is never negative.

Because the square-root sign gives a nonnegative answer, $\sqrt{x^2}$ is not always $x$: $\sqrt{(-3)^2}=\sqrt9=3$. In general, $\sqrt{x^2}=x$ when $x\ge0$, and $\sqrt{x^2}=-x$ when $x<0$, because then $-x$ is nonnegative and $(-x)^2=x^2$. In both cases $\sqrt{x^2}$ is the distance from $x$ to $0$ on the number line, which Section~\ref{sec:magnitude} calls the \emph{absolute value} $|x|$.
